\documentclass[11pt]{article}
\usepackage[a4paper, margin=1in]{geometry}
\usepackage{amsmath, amssymb}
\usepackage{enumitem}
\usepackage{fancyhdr}
\usepackage{xcolor}

\pagestyle{fancy}
\fancyhf{}
\lhead{CSE 400: Fundamentals of Probability in Computing}
\rhead{Lecture Scribe}
\cfoot{\thepage}

\newcommand{\solution}{
	\vspace{0.5cm}
	\noindent\textbf{\textcolor{blue}{Solution:}}
	\vspace{0.2cm}
}

\title{
	\normalsize School of Engineering and Applied Science (SEAS), Ahmedabad University \\
	\vspace{0.2cm}
	\textbf{CSE 400: Fundamentals of Probability in Computing}\\
	\Large Lecture Scribe Submission
}

\author{}
\date{}

\begin{document}
\maketitle

\vspace{-2cm}
\begin{center}
	\begin{tabular}{ll}
		\textbf{Name:} & {Kush Kelaiya \hspace{2.5in}} \\ [1.5ex]
		\textbf{Enrollment No:} & {AU2440018 \hspace{2.5in}} \\ [1.5ex]
		\textbf{Email:} & {kush.k2@ahduni.edu.in \hspace{2.1in}} \\ [1.5ex]
		\textbf{Group: } & s1\_g1\_fin \\ [1.5ex]
		\textbf{Lecture:} & 11 \\ [1.5ex]
		\textbf{Date of Submission:} & {February 17\textsuperscript{th}, 2026 (11:59PM) \hspace{1.2in}}
	\end{tabular}
\end{center}

\hrule
\vspace{0.5cm}

\section{Transformation of a Random Variable}
\subsection{Problem Setting}
Let:
\begin{itemize}
	\item $X$ be a continuous random variable (CRV).
	\item The PDF $f_X(X)$ of $X$ is known \textit{a priori}.
	\item Define a new random variable: $$Y = g(X)$$
	\item Objective: Find $$f_Y(Y),\,\,F_Y(Y)$$
\end{itemize}

\subsection{Case 1: $g(\cdot) Monotonically increasing$}
Assume:
\begin{itemize}
	\item $Y = g(X)$
	\item $g(\cdot)$ is \textbf{monotonically increasing}
	\item Hence inverse exists: $$X = g^{-1}(Y)$$
\end{itemize}

\textbf{Step 1: CDF of $Y$}
$$F_Y(T) = Pr(Y\leq y)$$
Substitute $Y = g(X)$
$$F_Y(Y) = Pr(g(X) \leq y)$$
Since $g(\cdot)$ is monotonically increasing:
$$g(X) \leq y \Longleftrightarrow X \leq g^{-1}(y)$$
Therefore,
$$
F_Y(Y) = Pr(X \leq g^{-1}(y))
F_Y(Y) = F_X(g^{-1}(y))
$$
\textbf{Step 2: Differentiate to Obtain PDF}
$$
f_Y(Y) = \dfrac{d}{dy}F_Y(Y)$$$$
f_Y(Y) = \dfrac{d}{dy}\left[ F_X(g^{-1}(y)) \right]
$$
Using chain rule:
$$f_Y(Y) = f_x(g^{-1}(y))\cdot\dfrac{d}{dy}[g^{-1}(y)]$$
Equivalently: $$f_Y(Y) = f_X(X)\left|\dfrac{dx}{dy}\right|_{x=g^{-1}(y)}$$
For increasing transformation: $$\dfrac{dx}{dy} > 0$$
Thus, $$f_Y(Y) = f_X(g^{-1}(y))\cdot\dfrac{dx}{dy}$$
\subsection{Case 2: $g(\cdot) Monotonically Decreasing$}
Assume:
\begin{itemize}
	\item $g(\cdot)$ is monotonically decreasing.
\end{itemize}
\textbf{Step 1: CDF}
\begin{align*}
	F_Y(Y) &= Pr(Y \leq y) \\
	&= Pr(g(X) \leq y)
\end{align*}
Since $g(\cdot)$ is decreasing:
$$g(X) \leq y \Longleftrightarrow X \geq g^{-1}(y)$$
Therefore,
$$F_Y(Y) = Pr(X \geq g^{-1}(y))$$
$$ = 1 - F_X(g^{-1}(y))$$
\textbf{Step 2: Differentiate}
\begin{align*}
	f_Y(Y) &= \dfrac{d}{dy}\left[ 1 - F_X(g^{-1}(y)) \right]\\
	&= -f_X(g^{-1}(y))\cdot\dfrac{d}{dy}[g^{-1}(y)]
\end{align*}
Since for decreasing transformation:
$$\dfrac{dx}{dy}<0$$
We write:
$$f_Y(Y) = f_X(g^{-1}(y))\left| \dfrac{dx}{dy} \right|$$

\subsection{General Formula}
\begin{center}
	\boxed{$$f_Y(Y) = f_X(g^{-1}(y))\left| \dfrac{dx}{dy} \right|$$}
\end{center}
with support determined by transformation limits.

\section{Worked Example}
\subsection{Given}
\begin{align*}
	X &\sim \text{Uniform}(-1, 1)\\
	f_X(X) &= \begin{cases}
		\dfrac{1}{2}, & -1 < x < 1 \\
		0, otherwise
	\end{cases}
\end{align*}
Define:
$$Y = sin\left( \dfrac{\pi X}{2} \right)$$
\subsection{Inverse Transformation}
\begin{align*}
	y = sin\left( \dfrac{\pi X}{2} \right) \\
	x = \dfrac{2}{\pi}sin^{-1}(y)
\end{align*}
\subsection{Compute Derivative}
$$
\dfrac{dx}{dy} = \dfrac{2}{\pi}\cdot\dfrac{1}{\sqrt{1 - y^2}}
$$
\subsection{Apply Formula}
$$
f_Y(Y) = f_X(X)\left| \dfrac{dx}{dy} \right|
$$
Since $fx(X) = \frac{1}{2}$ in support:
$$
f_Y(Y) = \dfrac{1}{2}\dfrac{2}{\pi\dfrac{1}{\sqrt{1 - y^2}}}
$$
\begin{center}
	\boxed{$$f_Y(Y) = \dfrac{1}{\pi\sqrt{1 - y^2}}$$}\;\;\; $-1 < y < 1$
\end{center}
Otherwise 0.
\section{Function of Two Random Variables}
\subsection{Definition}
Let:
$$Z = X + Y$$
Find:
\begin{enumerate}
	\item $f_Z(Z)$
	\item $F_Z(Z)$, if X and Y are independent
	\item if $X \sim N(0, 1)$, $Y \sim N(0, 1)$, independent, prove:
	$$Z \sim N(0, 2)$$
	\item if $X$, $Y$ exponential with parameter $\lambda$, find $f_Z(Z)$
\end{enumerate}
\subsection{Derivation Using CDF}
$$F_Z(Z) = Pr(Z \leq z) = Pr (X + Y \leq z)$$
From the geometric interpretation:\\
Region:
$$x + y \leq z$$
Therefore:
$$
F_Z(Z) = \int_{-\infty}^{\infty}\int_{-\infty}^{z-y}f_{XY}(x, y)dx\,dy
$$
Equivalently changing order:
$$
F_Z(Z) = \int_{-\infty}^{\infty}\int_{-\infty}^{z-x}f_{xy}(x, y)dy\,dx
$$
\subsection{PDF by Differentiation}
$$f_Z(Z) = \dfrac{d}{dx}F_Z(Z)$$
When $X$ and $Y$ are independent:
$$f_{XY}(x, y) = f_X(x)f_Y(y)$$
Hence convolution form:
\begin{center}
	\boxed{$$f_Z(Z) = \int_{-\infty}^{\infty}f_X(x)f_Y(z-x)dx$$}
\end{center}
\subsection{Normal Case}
Given:
$$
X \sim N(0, 1),\;\;\;Y \sim N(0, 1)
$$
Independent.
Using convolution:
$$
Z = X + Y
$$
Result:
\begin{center}
	\boxed{$$Z \sim N(0, 2)$$}
\end{center}
\subsection{Exponential Case}
If $X$, $Y$ exponential with parameter $\lambda$:
$$
f_Z(Z) = \int_{0}^{z}\lambda\epsilon^{-\lambda x}\epsilon^{-\lambda(z-x)}dx
$$

\vfill
\bigskip
\hrule
\vspace{0.2cm}
\begin{center}
\small \textit{End of Submission}
\end{center}

\end{document}